\subsubsection{Pentadic representation of the microfiber and its incidence}
\label{pi-estrella:desarrollo}

The arrangement of the five regions makes it possible to visualize simultaneously
the unity of the ternary quotient and the differentiation preserved by the
catalog. The relevant map is the restriction
\[
 \Psi\big|_{\mathscr R_\pi^{\rm micro}}:
 \mathscr R_\pi^{\rm micro}\longrightarrow\{010211\},
 \qquad \Psi(U)=(-U)\bmod3.
\]
Its five antecedents are the vectors \(U_6\) in the preceding table.
Figure \ref{pi-estrella:figura} recovers their radial organization,
preserving chart order, signatures, and multiplicities of the
catalog.

% Chart order and data recovered from fig_cinco_regiones_pi.tex.
\begin{figure}[H]
\centering
\begin{tikzpicture}[x=1cm,y=1cm,font=\small,>=Latex]
 \coordinate (a) at (90:2.2);
 \coordinate (b) at (18:2.2);
 \coordinate (c) at (-54:2.2);
 \coordinate (d) at (-126:2.2);
 \coordinate (e) at (162:2.2);
 \draw[black!35,line width=.45pt]
    (a)--(c)--(e)--(b)--(d)--cycle;
 \node[draw,fill=white,circle,minimum size=1.6cm,align=center,
    inner sep=3pt] (z) at (0,0) {$w_\pi$\\$010211$};
 \foreach \v in {a,b,c,d,e}{
    \fill (\v) circle (2pt);
    \draw[->,line width=.6pt] (\v)--(z);
 }
 \node[above=3mm,align=center] at (a)
   {$U_\pi^\perp$\\$555555$\\multiplicity $432$};
 \node[right=3mm,align=left] at (b)
   {$U_{\pi,1}^{++}$\\$339555$\\multiplicity $144$};
 \node[below right=2mm and 1mm,align=left] at (c)
   {$U_{\pi,2}^{++}$\\$393555$\\multiplicity $144$};
 \node[below left=2mm and 1mm,align=right] at (d)
   {$U_{\pi,3}^{++}$\\$933555$\\multiplicity $144$};
 \node[left=3mm,align=right] at (e)
   {$U_\pi^{--}$\\$933717$\\multiplicity $144$};
 \node[align=center,font=\footnotesize] at (0,-3.55)
   {$1_\perp+3_{++}+1_{--}$\qquad $432+4\cdot144=1008$};
\end{tikzpicture}
\caption{The five regions of \(\pi\) in a pentadic arrangement.
Each position preserves its multiplicative signature and multiplicity;
the radial arrows represent the common projection
\(\Psi(U)=010211\). The star-shaped outline organizes the five
positions of the chart. Their realization in the five Witt octads
is performed through regional alignment and the
\(4+1\) decomposition in the text.}
\label{pi-estrella:figura}
\end{figure}


Projection onto the center identifies the five ternary images; the
record of exterior positions retains the regional information
making incidence realization possible. The distribution
\(1_\perp+3_{++}+1_{--}\) specifies one transversal region, three
positive regions, and one negative region. Their normalized weights are
\(3/7\) and \(1/7\) for each of the remaining four: transversal
multiplicity is triple, even though all five images under
\(\Psi\) coincide.

\Needspace{11\baselineskip}
The link with the Witt star is formulated through the flag
\(B_K\subset H^-_\alpha\), the alignment \(\ell\), and the
lift \(\iota_D\) already constructed. The five regional positions
are realized in the five octads over \(\ell(B_K)\):
\[
\begin{aligned}
 U_\pi^\perp&\longmapsto O^\perp_{\ell(B_K)},\\
 U_\pi^{--}&\longmapsto\iota_D(H^-_\alpha),\\
 \{U_{\pi,1}^{++},U_{\pi,2}^{++},U_{\pi,3}^{++}\}
 &\longmapsto
 \{\iota_D(B_K\cup P):P\in\mathcal P_+\},\\
 \mathcal P_+&=\bigl\{\{1,3\},\{2,11\},\{8,12\}\bigr\}.
\end{aligned}
\]
Assignment of the three positive positions to the three pairs is
performed in an ordered chart. The transversal and negative rows
are distinguished by their regional predicates and by the flag.

The two figures are thus articulated: Figure
\ref{exc:fig-estrella} shows the four hexads of the small Witt
design over \(B_K\); Figure \ref{pi-estrella:figura} shows
the five regions that, under alignment, occupy the four residual
octads and the transversal octad of the binary design. The law
\(4+1\), proved in Proposition
\ref{exc:residuo-binario}, specifies the change in incidence. In this
composition, \(K\) determines the marked tetrad, and the signature of
\(\alpha\) distinguishes one of its hexads. In addition to their common
ternary value, the regions preserve the structure required for that
transport toward the exceptional realization.
